\chapter{Acrónimos y citas}

\section{Acrónimos}

Los acrónimos se gestionan mediante el paquete \texttt{acro}.

Por ejemplo, el acrónimo \ac{api} se expande automáticamente la primera vez
que aparece en el texto.

Esto sucede porque se inserta el código:
\begin{lstlisting}[language=TeX,caption={Ejemplo nuevo acrónimo}]
\DeclareAcronym{api}{
    short = API ,
    long  = Application Programming Interface 
}
\end{lstlisting} 

en la carpeta: 
\begin{lstlisting}[language=TeX,caption={Carpeta Acrónimos}]
\08_acronyms 
\end{lstlisting}

\section{Citas bibliográficas}

Las referencias bibliográficas se gestionan mediante \texttt{biblatex}.

Ejemplo de cita narrativa:

\begin{lstlisting}[language=TeX,caption={Cita narrativa}]
\textcite{knuth1984}
\end{lstlisting}

Ejemplo de cita parentética:

\begin{lstlisting}[language=TeX,caption={Cita parentética}]
\parencite{knuth1984}
\end{lstlisting}